\section{Related Work}
\label{sec:related}

\paragraph{C-MAPSS and RUL backbones.}
The C-MAPSS turbofan benchmark \cite{saxena2008} provides run-to-failure
trajectories for four subsets with different operating conditions and fault
modes. Most recent work trains a sequence model on the sensor trajectories and
predicts a capped RUL target. Reported gains are usually measured on a single
test window per engine, and the aggregation of multiple overlapping windows is
often treated as an implementation detail. Our protocol keeps the backbone frozen
and varies only the inference-time aggregation, so the comparison is not
confounded by backbone training.

\paragraph{RUL with uncertainty and domain shift.}
Evidential and domain-adaptation approaches model uncertainty or incomplete
degradation \cite{hou2025}; quantile and state-space approaches produce
distributional RUL estimates \cite{frizzo2025}; scientific machine learning
combines physics and data for engine health management \cite{barry2026}; and
survival-based aging models link degradation to RUL \cite{srinivasan2026}. These
lines improve the predictor itself. Our question is complementary: given a fixed
predictor, when does multi-window aggregation help, and when does a simple
post-hoc calibration suffice?

\paragraph{Aggregation and calibration.}
Fixed rules (mean, median, last value) and learned gates are common ways to
combine overlapping predictions. Post-hoc calibration (a linear or residual map
from predictions and context to the target) is equally common in deployment.
Calibration and reliability diagnostics have also been studied for RUL under
operating-regime shift \cite{yang2026}, as have conformal and similarity-based
interval methods \cite{busico2026,barber2023} and business-aware uncertainty
estimates \cite{frizzo2026}. These lines calibrate intervals or reliability;
they do not compare point-aggregation rankings across cuts.
The two families are rarely compared under the same cut protocol. We compare
them on the same frozen features, with the same engine-disjoint folds, and
report where each family wins. The result is a cut-dependent ranking rather than
a single winner.

\paragraph{Sensor-fault reconstruction.}
Measurement-model reconstruction uses a low-rank model of healthy sensor
correlations to replace a flagged sensor. It is attractive because it is
interpretable and can run online. The risk is that a temporal hold or median
baseline can recover the same signal without a model. We therefore compare the
reconstruction against fair backtracked baselines that use only observed data,
not against a weak last-value baseline. This is the comparison that removes the
apparent advantage on bias and drift.
